\chapter{Troubleshooting}
\label{ch:troubleshooting}

\section{The plugin does not find the server}

\begin{itemize}
  \item Is the server running? On its machine, \cmd{fxmenet-top} should show
    it. If not, start \cmd{fxmenet-server}.
  \item Does \ui{Find} show nothing? Discovery uses a broadcast, which does not
    cross routers or most VPNs, and is blocked by some firewalls. Type the
    server's address in the field instead.
  \item Can the client reach the server's ports? Check the firewall on the
    server (TCP 54322, UDP 54321 and 54320) and, on the client, that the
    DAW is allowed to use the network.
  \item Is a second server running on the same machine? Only one can answer
    \ui{Find}; give the other one a different port set, or no discovery.
\end{itemize}

\section{The plugin I want is not in the list}

\begin{itemize}
  \item It must be a VST3 \emph{effect} installed on the server, in a folder
    the server searches (section~\ref{sec:folders}).
  \item Look at the server's log (or the \key{3} view of
    \cmd{fxmenet-top}) for ``skipped'': a plugin that crashed or hung when
    loaded is left out until it changes.
  \item You installed it while the server ran: rescan with \key{r} in
    \cmd{fxmenet-top}, then reconnect the plugin.
  \item The list is filtered: clear the \ui{Filter} field.
\end{itemize}

\section{A plugin shows tier \tier{0} when I expected \tier{1}}

The plugin must be installed on the client, in a folder the client scans, in
the same version. Hover over the row for the reason. Use \ui{Rescan} or
\ui{Folders\ldots} in the \ui{Settings} tab (chapter~\ref{ch:tiers}).

\section{Concealed blocks, dropouts or clicks}

Look at where the problem is.

\begin{enumerate}
  \item In \cmd{fxmenet-top}, is the server's CPU high, or a session's load bar
    red? Then the server is overloaded: use larger blocks (1024$\times$2),
    reduce the load, or use a faster server.
  \item Does a session's \emph{lost} count rise? Requests do not reach the
    server in time: look at the network (wired? other traffic? a switch
    that drops packets?).
  \item Does \emph{send errors} rise in the server's log? The kernel drops
    replies: raise the socket buffers (section~\ref{sec:linuxtuning}).
  \item Do the threads show \cmd{TS} instead of \cmd{RR} in the check of
    section~\ref{sec:linuxtuning}? Give the audio threads real-time priority.
  \item Only the client's concealed count rises while the server's counts are
    zero? Replies are delayed on their way back, or the server stalls
    briefly. Raise D by one.
\end{enumerate}

If clicks persist, enable \ui{Log transport events} in the \ui{Settings}
tab, reproduce the problem, and keep \file{fxmenet-client-events.log} and the
server's log together for the report.

\section{Silence at start, stop or jump}

This is the \ui{Silence in-flight audio at reset} option of the
\ui{Settings} tab doing its job: for the length of the latency after a
restart, nothing is heard from the plugin, as with a local plugin whose delay
line is cleared. Turn the option off to hear what was still in flight.

\section{The audio is late compared with other tracks}

The DAW compensates for the latency that \fxme\ reports. If a track is
still late, check that the DAW's plugin delay compensation is on and that
the DAW applies it to this track.

\section{Everything stops when the server restarts}

The plugin reconnects by itself and rebuilds the chain on the server. If the
server does not have the same plugins any more, the affected plugins show an
error in the chain. While the server is away, the sound becomes dry at once; when the server is back, the processed sound returns. Restarting the server during playback was tested, without audible glitches.
